\scanpage{110}

Récapitulons le formulaire pour

\diagram{\begin{tikzcd}[ampersand replacement=\&,column sep=large]
k\arrow[r,"{\deleted{\unsure{a}}\alpha}"] \& A\ (\text{$k$-alg})\arrow[l,bend left=17,"t"]\arrow[l,bend left=42,"n"]
\end{tikzcd}}

\margin{$\alpha$ hom d'anneaux comm. $\alpha(\lambda)=\lambda.1_A$ On pose $\tau(x)=\alpha t(x)$ \ $\nu(x)=\alpha n(x)$}

\textcircled{1} $t$ \deleted{\unsure{est|$k$}} $k$-linéaire, $\boxed{t(1_A)=2}$

\textcircled{2} $\deleted{\unsure{\xi}}x\longmapsto\bar x\overset{\mathrm{def}}{=}\tau(x)-x:A\longrightarrow A$ est un automorphisme d'alg., i.e. on a

\margin{(autom. involutif du $k$-module $A$, avec $\bar1_A=1_A$)}

\[
\boxed{\overline{xy}=\bar x\bar y}\qquad\text{i.e.}\quad
2xy=(\tau(xy)-\tau(x)\tau(y))+\tau(y)x+\tau(x)y
\]

\[
\text{\textcircled{3}}\quad\boxed{\left\{\begin{aligned}
n(\lambda x)&=\lambda^2n(x)\\
n(x+y)&=n(x)+n(y)+t(x\bar y)
\end{aligned}\right.}\qquad\lambda\in k,\ x\in A
\]

i.e. $n$ est quadratique, la forme bil. ass. est ici

\[
\varphi_n(x,y)=t(\unsure{x\bar y|y\bar x})
\]

\textcircled{4} $\boxed{n(1)=1}$ (\uline{NB} compte \deleted{\unsure{tenu}} de la formule d'additivité du 3°, cela implique $t(1)=2$)

\textcircled{5} $\boxed{\alpha(n(x))=x\bar x}$ $(x\in A)$ qui implique 3°, 4° si $\alpha$ injectif

\textcircled{6} (\uline{facultatif}, \deleted{)}\inserted{et} si $k$ loc. libre de rang $n$ sur $k_0$, et $A$ loc. libre de rang $2n$ sur $k_0$)

\[
\boxed{\begin{aligned}
\mathrm{Tr}_{k/k_0}(t(x))&=\mathrm{Tr}_{A/k_0}(x)\\
N_{k/k_0}(n(x))&=N_{A/k_0}(x)
\end{aligned}}\qquad x\in A
\]

La donnée de \inserted{\unsure{$\uparrow$} $A$ sur $k$ avec} $\alpha,t,n$ satisfaisant 1° à 5° sera appelée algèbre \uline{quasi-quadratique} sur $k$.

On va maintenant pouvoir étudier la situation \inserted{quasi-}\uline{biquadratique}

% archive sheet number: 109

